% This document was generated by an LLM.
% Content fragment; compile intro-analysis-self-check.tex.
% Coverage: Hunter (2012), Chapter 4, printed pp. 39--55.
% Corrections and clarifications:
% - In Example 4.14 the extension is discontinuous at 1, so its right
%   derivative there does not exist; the printed value 0 is incorrect.
% - The product-rule proof uses Theorem 4.16 for continuity, not 4.18.
% - In the chain-rule proof, t(h)=g'(f(c))r(h)+s(phi(h)); the printed
%   expression omits the coefficient of r(h). Also handle phi(h)=0.
% - Taylor polynomials have degree at most n, not necessarily exactly n.
% - In the Taylor proof assume x != c; the case x=c is immediate.

\section{Differentiable Functions}

\subsection{The derivative}
\textit{Source: \S4.1, pp.~39--45.}

\subsubsection*{Definition and examples}
\begin{enumerate}
  % Definition 4.1.
  \item Define differentiability of \(f:(a,b)\to\R\) at \(c\),
    the derivative \(f'(c)\), and differentiability on \((a,b)\).
    What is the domain of the derivative function?

  % Equivalent formula and geometric discussion following Definition 4.1.
  \item Give the two equivalent difference-quotient definitions of
    \(f'(c)\), using \(h\to0\) and \(x\to c\).
    Why need neither quotient be defined at its limiting input?
    Explain the secant-slope and tangent-slope interpretation.

  % Locality discussion on pp. 39--40.
  \item What does it mean that differentiability is a local property?
    Explain how the derivative can be defined at any interior point
    of a domain \(A\subseteq\R\), even when \(A\) is not an interval.

  % Example 4.2.
  \item Differentiate \(f(x)=x^2\) directly from the definition.
    Why is cancelling a factor of \(h\) legitimate before evaluating
    the limit, while simply substituting \(h=0\) into the original
    quotient is not?

  % Example 4.3.
  \item Let \(f(x)=x^2\) for \(x>0\) and \(f(x)=0\) for
    \(x\leq0\). Find its derivative at every real point, using
    one-sided limits of the difference quotient at zero.

  % Example 4.4.
  \item Let \(f(0)=0\) and \(f(x)=1/x\) for \(x\ne0\).
    Determine where it is differentiable and compute the derivative
    directly from difference quotients. What happens to the quotient
    at zero?

  % Examples 4.5--4.6.
  \item Determine the differentiability points and derivatives of
    \(\operatorname{sgn}x\) and \(|x|\). Compare their continuity
    and difference quotients at zero.

  % Example 4.7.
  \item Differentiate \(x^{1/3}\) at nonzero points using the
    difference-of-cubes identity. What happens to the difference
    quotient at zero? Does an infinite limiting slope count as a
    derivative in the definition used here?

  % Example 4.8.
  \item Let \(f(0)=0\) and \(f(x)=x\sin(1/x)\) for \(x\ne0\).
    Determine its differentiability points and its derivative where
    defined. Prove the conclusion at zero from the definition;
    the product and chain rules may be used away from zero.

  % Example 4.9.
  \item Let \(f(0)=0\) and \(f(x)=x^2\sin(1/x)\) for \(x\ne0\).
    Compute \(f'\) at all real points. Is \(f'\) continuous at zero?
    Prove both the existence of \(f'(0)\) and your conclusion about
    continuity of the derivative.
\end{enumerate}

\subsubsection*{Derivatives as linear approximations}
\begin{enumerate}[resume]
  % Definition of little-o preceding Proposition 4.10.
  \item Define \(r(h)=o(h)\) as \(h\to0\).
    Is \(r(h)\to0\) alone sufficient to imply \(r(h)=o(h)\)?
    Give an example that distinguishes the conditions.

  % Proposition 4.10.
  \item State the characterization of differentiability by a linear
    approximation and a remainder (Proposition 4.10).
    What is the coefficient of the linear term, and what condition
    must the remainder satisfy?

  % Proof of Proposition 4.10.
  \item Prove both directions of Proposition 4.10.
    How does dividing the error by \(h\) connect the linear
    approximation to the difference-quotient definition?

  % Geometric and analytic discussion on p. 43.
  \item Explain how \(f'(c)\) describes the first-order change in
    \(f(c+h)-f(c)\). What do its magnitude and sign say about
    infinitesimal scaling and orientation? Distinguish the affine
    approximation to \(f(c+h)\) from the linear approximation to
    its change.

  % Examples 4.11--4.12.
  \item For \(f(x)=x^2\), and then for \(f(x)=1/x\) at \(c\ne0\),
    write \(f(c+h)=f(c)+f'(c)h+r(h)\) with an exact expression
    for \(r(h)\). Verify the required little-oh condition.
\end{enumerate}

\subsubsection*{Left and right derivatives}
\begin{enumerate}[resume]
  % Definition 4.13.
  \item Define the right and left derivatives of a function on
    \([a,b]\), including where each can be considered.
    Write the left derivative both as a limit with \(h\to0^-\)
    and as a limit with \(h\to0^+\).

  % Observation following Definition 4.13 and Example 4.15.
  \item Give a necessary and sufficient condition for a two-sided
    derivative to exist in terms of one-sided derivatives.
    Illustrate failure of the condition with \(|x|\) at zero.

  % Example 4.14, correcting the right derivative at 1 for its extension.
  \item Find the endpoint derivatives of \(f(x)=x^2\) on \([0,1]\).
    Now extend it by setting \(f(x)=1/x\) for \(x<0\) and
    \(f(x)=0\) for \(x>1\). Which one-sided derivatives at zero
    and one still exist, and what are their values?
    Check the right-derivative claim at one in Example 4.14 using
    the actual value \(f(1)\).
\end{enumerate}

\subsection{Properties of the derivative}
\textit{Source: \S4.2, pp.~45--49.}

\subsubsection*{Differentiability and continuity}
\begin{enumerate}[resume]
  % Theorem 4.16 and its proof.
  \item State and prove the implication from differentiability to
    continuity at a point (Theorem 4.16). How can the increment
    \(f(c+h)-f(c)\) be written using the difference quotient?

  % Counterexamples following Theorem 4.16.
  \item Does continuity at a point imply differentiability there?
    Give examples from this chapter exhibiting different ways that
    differentiability can fail for a continuous function.

  % Definition 4.17.
  \item Define continuous differentiability and the notation
    \(f\in C^1(a,b)\). Give an example of a differentiable
    function that is not continuously differentiable.
\end{enumerate}

\subsubsection*{Algebraic properties}
\begin{enumerate}[resume]
  % Theorem 4.18.
  \item State the scalar-multiple, sum, product, and quotient rules
    for derivatives (Theorem 4.18), with all hypotheses.

  % Product-rule proof in Theorem 4.18.
  \item Prove the product rule from the difference-quotient definition.
    Which term do you add and subtract, and where is continuity of
    a differentiable function needed?

  % Example 4.19.
  \item Derive the power rule for positive integer powers by induction
    using the product rule. Deduce differentiability of polynomials
    and rational functions on their natural domains.
\end{enumerate}

\subsubsection*{The chain rule}
\begin{enumerate}[resume]
  % Theorem 4.20.
  \item State the chain rule (Theorem 4.20), including the domain,
    interior-point, and differentiability assumptions.
    At which point is the derivative of the outer function evaluated?

  % Linear-approximation motivation preceding Theorem 4.20.
  \item Explain the chain-rule formula using composition of the
    linear approximations to the inner and outer functions.

  % Proof of Theorem 4.20; correct the omitted coefficient.
  \Needspace{9\baselineskip}
  \item Prove the chain rule using
    \[
      f(c+h)=f(c)+f'(c)h+r(h),\qquad
      g(f(c)+k)=g(f(c))+g'(f(c))k+s(k).
    \]
    Identify the full remainder after composition, including the
    coefficient of \(r(h)\). Show that it is \(o(h)\).
    Your argument should also cover \(f'(c)=0\) and increments
    for which \(f(c+h)=f(c)\).

  % Example 4.21 and the quotient-rule discussion.
  \item Use the chain rule to differentiate \(1/f\) where \(f\)
    is differentiable and nonzero. Deduce the quotient rule from
    the reciprocal formula and the product rule.
\end{enumerate}

\subsubsection*{The derivative of inverse functions}
\begin{enumerate}[resume]
  % Proposition 4.22.
  \item State the formula for the derivative of an inverse function
    (Proposition 4.22), with every assumption in the proposition.
    Does this proposition itself establish differentiability of the inverse?

  % Proof and alternative notation following Proposition 4.22.
  \item Derive the inverse-derivative formula from the identity
    \(f^{-1}(f(x))=x\). Write the formula at \(d=f(c)\) as well
    as at \(f(c)\). Why is differentiability of the inverse impossible
    there if \(f'(c)=0\) under the stated interior-point assumptions?

  % Example 4.23.
  \item Apply the inverse-derivative formula to \(f(x)=x^3\).
    Where is its inverse differentiable, and what happens at zero?

  % First unproved result following Example 4.23.
  \item State the monotonicity and inverse-continuity facts for a
    continuous one-to-one function on an interval.
    Under these hypotheses, what extra pointwise conditions guarantee
    differentiability of the inverse at an interior image point?

  % One-dimensional inverse function theorem stated at the end of Section 4.2.
  \item State the local inverse function theorem quoted in the notes.
    Specify the differentiability and continuity assumptions near
    the point and the conclusion about the local inverse and its derivative.
\end{enumerate}

\subsection{Extreme values}
\textit{Source: \S4.3, pp.~49--51.}

\begin{enumerate}[resume]
  % Definition 4.24.
  \item Define global and local maxima and minima of \(f:A\to\R\).
    What is an extreme value? Distinguish the point where an extremum
    is attained from the extreme value itself.

  % Theorem 4.25.
  \item State Fermat's necessary condition for a local extremum
    (Theorem 4.25). Include the assumptions about the point and
    differentiability of the function.

  % Proof of Theorem 4.25.
  \item Prove Fermat's theorem using signs of the difference quotient
    on the two sides of a local maximum. How does the proof change
    for a local minimum?

  % Endpoint discussion following Theorem 4.25.
  \item Why does Fermat's zero-derivative conclusion require an
    interior point? State the necessary one-sided derivative
    inequalities at each endpoint of \([a,b]\) for a local maximum
    and for a local minimum, assuming the relevant derivatives exist.

  % Definition 4.26.
  \item Define critical and stationary points, using the convention
    in these notes. Must a critical point have a derivative?
    Are endpoints included in either definition?

  % Candidate list following Definition 4.26; test necessary vs sufficient.
  \item List the three types of points that must be considered when
    searching for attained extrema on a domain \(A\subseteq\R\).
    Does being stationary guarantee an extremum? Give a counterexample.
    Give an example of an interior extremum at a point where the
    function is not differentiable.
\end{enumerate}

\Needspace{8\baselineskip}
\subsection{The mean value theorem}
\textit{Source: \S4.4, pp.~51--53.}

\begin{enumerate}[resume]
  % Theorem 4.27.
  \item State Rolle's theorem (Theorem 4.27) on a nondegenerate
    interval \([a,b]\). Distinguish the continuity requirement at
    the endpoints from the differentiability requirement in the interior.

  % Proof of Theorem 4.27.
  \item Prove Rolle's theorem using the extreme value theorem and
    Fermat's theorem. Handle the constant-function case explicitly.
    Where does this argument ultimately depend on completeness of \(\R\)?

  % Theorem 4.28.
  \item State the mean value theorem (Theorem 4.28), including every
    hypothesis and the location of the point it supplies.
    Give its geometric interpretation.

  % Proof of Theorem 4.28.
  \item Deduce the mean value theorem from Rolle's theorem.
    Which affine function should be subtracted from \(f\), and
    why does the resulting function meet Rolle's hypotheses?

  % Analytic interpretation following Theorem 4.28.
  \item Explain how the mean value theorem connects local derivative
    information to the change between two function values.
    If \(|f'(x)|\leq M\) on \((a,b)\) and \(f\) is continuous
    on \([a,b]\), what bound does it give for \(|f(b)-f(a)|\)?

  % Theorem 4.29 and its proof.
  \item State and prove the zero-derivative theorem (Theorem 4.29).
    On what kind of domain is its constancy conclusion asserted?
    Give an example showing why a domain consisting of two disjoint
    open intervals can invalidate that conclusion.

  % Corollary 4.30 and its proof.
  \item State and prove the result for two differentiable functions
    with equal derivatives on an interval (Corollary 4.30).
    How does the zero-derivative theorem apply?

  % Theorem 4.31.
  \item State the derivative-sign characterizations of increasing
    and decreasing differentiable functions on an interval
    (Theorem 4.31). State also the sufficient conditions for strict
    increase and strict decrease.

  % Proof of Theorem 4.31.
  \item Prove both directions of the increasing-function
    characterization. Which direction follows from difference
    quotients, and which uses the mean value theorem?
    Explain the strict and decreasing versions.

  % Example 4.32.
  \item Must the derivative of a differentiable strictly increasing
    function be positive at every point? Give a counterexample and
    state the derivative inequality that does hold everywhere.

  % Discussion preceding Example 4.33.
  \item Suppose \(f\) is differentiable near \(c\), \(f'(c)>0\),
    and \(f'\) is continuous at \(c\). Prove that \(f\) is
    strictly increasing in some neighborhood of \(c\).
    Where is continuity of the derivative used?

  % Example 4.33.
  \item Let \(f(0)=0\) and
    \(f(x)=x/2+x^2\sin(1/x)\) for \(x\ne0\).
    Show that \(f'(0)>0\) but \(f\) is not increasing on any
    neighborhood of zero. Find nonzero points arbitrarily close to
    zero where the derivative is negative.
\end{enumerate}

\subsection{Taylor's theorem}
\textit{Source: \S4.5, pp.~53--55.}

\begin{enumerate}[resume]
  % Opening paragraph of Section 4.5.
  \item Define the second derivative, higher-order derivatives,
    and infinite differentiability on an open interval.

  % Definition 4.34; allow lower actual degree.
  \item Define the order-\(n\) Taylor polynomial of \(f\) at \(c\)
    and its Taylor coefficients. Give the summation formula.
    Must its actual polynomial degree equal \(n\), or can it be smaller?

  % Example 4.35.
  \item If \(P\) is a polynomial of degree at most \(n\), what
    is its order-\(n\) Taylor polynomial about an arbitrary center
    \(c\)? Verify the claim.

  % Example 4.36.
  \item Compute the order-\(n\) Taylor polynomial of \(e^x\)
    at zero from the definition of the coefficients.

  % Examples 4.37--4.38.
  \item Compute the Taylor polynomials of \(\cos x\) and \(\sin x\)
    at zero. Specify the nonzero coefficients and explain which
    successive orders give identical polynomials.

  % Example 4.39.
  \item Compute the order-\(n\) Taylor polynomial of \(1/x\) at
    \(c=1\), deriving the coefficients from its higher derivatives.

  % Example 4.40.
  \item Compute the order-\(n\) Taylor polynomial of \(\log x\)
    at \(c=1\), including its constant term and the general nonzero
    coefficient.

  % Remainder definition preceding Theorem 4.41.
  \item Define the Taylor remainder \(R_n(x)\).
    Why does calculating a Taylor polynomial alone not establish
    the accuracy of an approximation away from its center?

  % Theorem 4.41.
  \item State Taylor's theorem with Lagrange remainder
    (Theorem 4.41) for a function with \(n+1\) derivatives on an
    open interval. Specify the location of the intermediate point
    for \(x\ne c\). Does the statement require continuity of
    the \((n+1)\)st derivative?

  % Mean-value interpretation preceding Theorem 4.41.
  \item Show that the case \(n=0\) of Taylor's theorem is the mean
    value theorem applied between \(c\) and \(x\).

  % First part of the proof of Theorem 4.41.
  \item In the proof of Taylor's theorem, fix \(x\ne c\) and define
    \[
      g(t)=f(x)-\sum_{k=0}^{n}\frac{f^{(k)}(t)}{k!}(x-t)^k.
    \]
    Compute \(g(x)\), \(g(c)\), and \(g'(t)\).
    Explain the cancellation in the derivative and the relationship
    between \(g(c)\) and the desired remainder.

  % Remaining part of the proof of Theorem 4.41.
  \item Complete the proof by constructing from \(g(t)\) an auxiliary
    function that vanishes at \(c\) and \(x\), as in the notes.
    Apply Rolle's theorem and derive the Lagrange remainder formula.
    Why is cancellation of \((x-\xi)^n\) legitimate, and how is
    the case \(x=c\) handled separately?

  % Discussion of the Lagrange remainder on p. 55.
  \item Suppose \(|f^{(n+1)}(t)|\leq M\) between \(c\) and \(x\).
    Derive an explicit bound on \(|R_n(x)|\).
    How does the Lagrange remainder differ from simply using the
    next Taylor coefficient evaluated at \(c\)?

  % Example 4.42.
  \item Use Taylor's theorem to compute
    \[
      \lim_{x\to0}\frac{1-\cos x}{x^2}.
    \]
    Give an explicit error bound for the difference between the
    quotient and its limit, valid for nonzero \(x\).
    Which Taylor order gives the quadratic error estimate in the notes?
\end{enumerate}
